\section{宪法的实施及其保障}
\begin{enumerate}
  \item 修改、解释、监督、宣誓。
  \item 宪法解释。
        \begin{enumerate}
          \item 议会模式/立法机关模式。英国。
          \item 司法机关模式。美国。
          \item 专门机关模式。法国、德国。宪法法院。
          \item 归全人常。但没有具体的程序。
        \end{enumerate}
  \item 全人常的宪法监督程序。模式与解释相同。违宪审查。
        \begin{enumerate}
          \item 依据。宪法规定。
          \item 事前、事后审查。
          \item 下位法可以\zd{主动审查}。专项审查。
          \item 被动审查。
                \begin{enumerate}
                  \item 审查要求。两央、两高、省常。
                  \item 审查建议。
                  \item 有问题书面审查意见 \yjt 不改 \yjt 提案给委员长会议撤销（最后沟通）\yjt 常委会。
                  \item 应当反馈，可以公开。
                \end{enumerate}
          \item （法院）附带性审查和（公民）宪法控诉。
        \end{enumerate}
  \item 宪法宣誓。谁产生谁负责组织。
        \begin{enumerate}
          \item 国家机关工作人员；就职时；公开；法律规定。
          \item 驻外全权代表。外交部。
          \item 两院一委自己。
        \end{enumerate}
\end{enumerate}
